\chapter{Method}
\label{chap:method}

\section{Data Pipeline}

Preprocessing extracts block textures from base and target resource packs, filters to full-cube block textures, resizes images to \(32 \times 32\), and writes per-pack datasets.
Each training sample pairs one vanilla texture with the matching texture from one target pack.
The pipeline also records metadata for style-reference sampling and validation.

\section{Model Overview}

The current system uses two complementary models:

\begin{itemize}
  \item \textbf{StyleAwareUNet}: a diffusion model that fuses noisy target state with vanilla content, encodes multiple target-pack style references, and applies cross-attention at the bottleneck.
  \item \textbf{RecolorNet}: a lightweight U-Net that predicts a residual over the vanilla content texture and acts as a fallback when diffusion candidates are structurally weak.
\end{itemize}

\placeholderfigure{SpriteCraft generation overview}{Replace this placeholder with a pipeline diagram showing vanilla content input, style references, diffusion candidates, recolor fallback, and final candidate selection.}

\section{Conditioning}

The diffusion model receives the current noisy target image, the vanilla content image, three ranked style references from the same target pack, a validity mask for padded references, and the diffusion timestep.
Reference ranking uses content descriptors and filename-family priors so that visually relevant materials are more likely to guide generation.

\section{Training Objective}

The diffusion loss combines noise prediction with auxiliary reconstruction and content-aware penalties:

\begin{equation}
  \mathcal{L}_{total}
  =
  \mathcal{L}_{noise}
  + 0.70\mathcal{L}_{recon}
  + 0.30\mathcal{L}_{gradient}
  + 0.30\mathcal{L}_{content}.
\end{equation}

The content term includes structure, contrast, hue, and color-moment components.
The important design choice is that structure comparisons are source-relative, while contrast and hue penalties are target-gated.
This focuses the model on changes that matter for the target style without forcing flat or low-chroma regions to carry unnecessary detail.

\section{Inference}

Sampling starts from Gaussian noise and applies 40 reverse diffusion steps.
The system samples multiple diffusion candidates, scores them against style support descriptors, optionally compares against a learned recolor candidate, and falls back to deterministic recoloring when necessary.
The report should include qualitative examples whenever discussing this selection logic because scalar scores can miss tileability and material identity failures.
